% =====================================================================
%  Método 9 — Surpresa por token   (dimensão: semântica)
%  Origem: células 41, 42 e 43 do caderno
% =====================================================================
\subsection{Surpresa por token}
\label{met:surpresa}

O SLOR e o PLL devolvem um valor para a frase integral. Prestam-se à ordenação de frases, mas
não localizam o desvio. A surpresa devolve um valor por palavra: é a mesma grandeza,
examinada antes da agregação.

Para cada posição, a surpresa é o logaritmo negativo da probabilidade que o modelo atribui à
palavra efetivamente ocorrente:
%
\begin{equation}
\surp(w_i) = -\ln P\bigl(w_i \mid w_{i-1}\bigr)
\end{equation}
%
Como $P \in (0,1]$, o logaritmo é negativo e a surpresa é sempre não negativa. Assume valor
nulo sob certeza do modelo e cresce ilimitadamente à medida que a palavra se torna improvável.
A unidade é relevante: em logaritmo natural a grandeza é expressa em \emph{nats} e, sob
$\log_2$, em \emph{bits}, caso em que a leitura é literal --- a quantidade de informação que
a palavra veicula. A surpresa corresponde exatamente à autoinformação da teoria da informação.

O método não constitui técnica nova, mas a parcela que o SLOR já agregava. Retomando os dois
estimadores da subseção~\ref{met:slor},
%
\begin{equation}
\SLOR(\sent) = \frac{1}{|\sent|}\sum_{i} \Bigl[\,\surp_{\mathrm{uni}}(w_i)
- \surp_{\mathrm{bi}}(w_i)\,\Bigr]
\end{equation}
%
isto é, o SLOR é a média da redução de surpresa produzida pelo contexto. Este método é o mesmo
cálculo exposto antes da agregação, de modo que nada é recomputado e nenhum modelo adicional
se faz necessário.

A leitura proposta consiste em identificar o pico,
%
\begin{equation}
\hat{\imath} \;=\; \arg\max_{i} \; \surp(w_i)
\end{equation}
%
sob a expectativa de que a surpresa se eleve na posição defeituosa, o que permitiria localizar
a palavra. Convém registrar a suposição embutida nessa leitura: assume-se que \emph{improvável}
e \emph{incorreto} coincidam.

O método devolve um perfil com um valor por palavra, do qual se retém o pico. A
Tabela~\ref{tab:surpresa} apresenta o resultado sobre três frases.

\begin{table}[H]
\centering
\dimtable{
\begin{tabular}{@{}lrl@{}}
\toprule
\textbf{Frase} & \textbf{Pico} & \textbf{Palavra} \\
\midrule
Há três dias o incêndio atingiu uma área de mata nativa & 11,47 & \texttt{atingiu} \\
Há três dias o incêndio cantou uma área de mata nativa  & 13,70 & \texttt{cantou} \\
Focos recorrentes apresentaram persistência anômala     & 15,08 & \texttt{anômala} \\
\bottomrule
\end{tabular}}
\caption{O pico de surpresa em três frases, em nats.}
\label{tab:surpresa}
\end{table}

Nas posições em que o contexto efetivamente intervém, a leitura é adequada. Isoladamente,
\texttt{nativa} é palavra de frequência muito baixa, com surpresa de $11{,}58$; na sequência
de \texttt{mata}, contudo, é praticamente esperada e o valor decresce para $3{,}53$. É o
comportamento pretendido pela técnica, a saber, mensurar a palavra em posição, e não em
abstrato. As duas outras linhas da tabela, entretanto, correspondem a frases corretas, e a
Seção~\ref{sec:empiricos} examina o que o $\arg\max$ está de fato selecionando.
